\section{Capítulo II: Aspectos técnicos y arquitectura de la red}
En este segundo capítulo de la monografía se abordarán los aspectos técnicos y tecnológicos más relevantes de la Red Dorsal Nacional de Fibra Óptica (RDNFO), con el propósito de comprender cuáles fueron los retos durante su desarrollo y bajo qué objetivos de conectividad se diseñó este proyecto. Para ello, primero analizaremos los componentes físicos de la red y el uso de la tecnología de multiplexación DWDM para maximizar su capacidad de transmisión. Luego, se examinará la arquitectura y cobertura del sistema, detallando su topología en anillos y cómo se planificó su interconexión con las redes regionales a nivel nacional. Finalmente, abordaremos el modelo de operación mayorista bajo el rol de "operador de operadores" y cómo la estructura de tarifa única influyó en su viabilidad dentro del mercado de telecomunicaciones.
\subsection{Componentes tecnológicos y medio físico}
\subsubsection{Características de la fibra óptica y métodos de tendido}
Para garantizar la viabilidad y resistencia de la RDNFO a lo largo del accidentado territorio peruano, se optó por la instalación de un cable de fibra óptica monomodo de 48 hilos. Específicamente, este medio físico cumple con los estándares internacionales de las recomendaciones ITU-T G.652.D y G.655, tecnologías diseñadas para transmisiones de larga distancia y alta capacidad. Además, según detalla Chávez (2016), la red dorsal emplea cable dieléctrico ADSS (dieléctrico autosoportado) de núcleo seco con gel de relleno, lo cual le otorga protección contra la humedad y evita cualquier tipo de inducción eléctrica al no poseer elementos metálicos. 

Respecto al método de tendido, uno de los mayores retos del proyecto fue el despliegue geográfico. Para superar esto de forma eficiente, se implementó un tendido aéreo aprovechando la infraestructura ya existente en el país. De este modo, la fibra se instaló sobre las torres de alta y media tensión pertenecientes a las empresas distribuidoras de electricidad. Este mecanismo de "compartición de infraestructura" estuvo respaldado por el marco normativo de la Ley N° 29904, permitiendo una instalación más rápida y económica en comparación con un tendido subterráneo tradicional.
\subsubsection{Tecnología de multiplexación DWDM y capacidad de transmisión}
La multiplexación por división densa de longitud de onda (Dense Wavelength Division Multiplexing, DWDM) es una tecnología de transmisión óptica que incrementa significativamente la capacidad de las redes de fibra óptica al permitir la transmisión simultánea de múltiples señales sobre una misma fibra mediante diferentes longitudes de onda. Cada longitud de onda funciona como un canal independiente, lo que posibilita aprovechar al máximo la infraestructura instalada sin necesidad de desplegar nuevos enlaces físicos. Para ello, la arquitectura DWDM emplea componentes como multiplexores, demultiplexores, transpondedores y amplificadores ópticos de fibra dopada con erbio (EDFA), los cuales permiten combinar, separar, convertir y amplificar las señales ópticas, garantizando transmisiones de alta velocidad y larga distancia con elevados niveles de eficiencia y confiabilidad (Agrawal, 2012; Keiser, 2011).
Además, una de las principales ventajas de esta tecnología es su escalabilidad, ya que la capacidad del enlace puede incrementarse incorporando un mayor número de longitudes de onda sobre la misma fibra, sin realizar modificaciones en la infraestructura física. De esta manera, los sistemas DWDM pueden alcanzar capacidades de transmisión de varios terabits por segundo, razón por la cual se han convertido en la tecnología predominante para las redes troncales (backbone) de operadores y proveedores de servicios de telecomunicaciones a nivel mundial (Keiser, 2011; International Telecommunication Union [ITU-T], 2020).
En adición a ello, la capacidad de transmisión se define como la cantidad máxima de información que un sistema de comunicaciones puede transportar a través de un medio físico en un intervalo determinado de tiempo. En redes de fibra óptica, esta capacidad se expresa en bits por segundo (bps) y depende de las características del medio de transmisión y de las tecnologías empleadas, como la multiplexación por división densa de longitud de onda (DWDM), que permite incrementar el ancho de banda al transmitir múltiples canales ópticos sobre una misma fibra (Keiser, 2011; Agrawal, 2012).
\subsection{Arquitectura de red y cobertura geográfica}
\subsubsection{Topología en anillos, nodos y redundancia del sistema}
La arquitectura de una red de transporte de fibra óptica comprende la organización física y lógica de sus enlaces, nodos y mecanismos de protección, con el propósito de garantizar la continuidad del servicio y optimizar el transporte de información. En redes troncales de gran escala, como la Red Dorsal Nacional de Fibra Óptica (RDNFO), la topología empleada responde a criterios de disponibilidad, escalabilidad y tolerancia a fallas, características indispensables para soportar el tráfico nacional de telecomunicaciones. En este contexto, la infraestructura se compone de nodos de agregación y transporte conectados mediante enlaces de fibra óptica de alta capacidad que conforman la red backbone del país (Pacheco et al., 2017). 
La RDNFO fue diseñada bajo una topología en anillos, configuración ampliamente utilizada en redes de transporte debido a que proporciona redundancia y continuidad operativa. En este tipo de arquitectura, los nodos se encuentran conectados mediante rutas alternativas que permiten redirigir automáticamente el tráfico cuando ocurre una interrupción en alguno de los enlaces. De esta manera, ante un corte de fibra óptica o una falla en un equipo de transmisión, la información puede circular por el sentido opuesto del anillo, reduciendo significativamente la probabilidad de interrupciones del servicio y aumentando la disponibilidad de la red (Ministerio de Transportes y Comunicaciones [MTC], 2013). 
Dentro de esta arquitectura, los nodos constituyen los puntos donde convergen, procesan y distribuyen los flujos de información provenientes de distintas rutas de transmisión. Estos nodos incorporan tecnologías de transporte óptico y conmutación, como DWDM y Metro Ethernet, permitiendo la agregación del tráfico regional y su encaminamiento hacia otras zonas del país. Según OSIPTEL, la RDNFO cuenta con 322 nodos distribuidos estratégicamente para interconectar las 180 capitales de provincia, conformando la infraestructura nacional de transporte sobre la cual operan diversos proveedores de servicios de telecomunicaciones (Pacheco et al., 2017). 
La confiabilidad del sistema también depende de sus mecanismos de redundancia, entendidos como el conjunto de recursos físicos y lógicos que garantizan la continuidad del servicio frente a eventos imprevistos. En la RDNFO, la existencia de rutas alternativas, enlaces redundantes y una disponibilidad mínima anual del 99,999 % para la red principal permiten mantener la operación incluso cuando se presentan fallas en determinados segmentos. Asimismo, el contrato de concesión estableció que, cuando la utilización de un enlace alcanzara el 75 % de su capacidad instalada, el operador debía ampliar su capacidad para preservar los niveles de desempeño y evitar situaciones de congestión (MTC, 2013). 

\subsubsection{Despliegue nacional e interconexión con las Redes Regionales}
El despliegue de la Red Dorsal Nacional de Fibra Óptica representa uno de los proyectos de infraestructura de telecomunicaciones más importantes desarrollados en el Perú, al constituirse como la red de transporte encargada de interconectar las principales ciudades del territorio nacional. La RDNFO dispone de 13 571 km de fibra óptica intermetropolitana, conectando 23 departamentos, 180 provincias y 284 distritos mediante una infraestructura que enlaza todas las capitales provinciales del país. Esta red fue concebida para proporcionar capacidad mayorista a los operadores de telecomunicaciones y servir como soporte para la expansión de los servicios de banda ancha en zonas urbanas y rurales (Pacheco et al., 2017). 
No obstante, la cobertura nacional no depende únicamente de la Red Dorsal. Con el objetivo de extender la conectividad hacia localidades con menor acceso a infraestructura de telecomunicaciones, el Estado peruano impulsó el desarrollo de los Proyectos Regionales de Banda Ancha, cuya finalidad consiste en ampliar la red de fibra óptica desde las capitales provinciales hasta las capitales distritales y otros centros poblados. Estos proyectos contemplaron el despliegue de más de 28 000 km adicionales de fibra óptica, permitiendo incrementar la capilaridad de la red y brindar cobertura a aproximadamente 1 600 distritos de los 1 834 existentes en el país (MTC, 2013). 
La interconexión entre la RDNFO y las Redes Regionales responde a un modelo jerárquico de transporte, donde la red dorsal actúa como backbone nacional mientras que las redes regionales distribuyen la conectividad hacia ámbitos de menor cobertura. Esta integración facilita el transporte eficiente del tráfico de datos entre las distintas regiones, fortalece la infraestructura nacional de telecomunicaciones y permite que los operadores locales accedan a una red de transporte de alta capacidad para ofrecer servicios de Internet, telefonía y transmisión de datos a usuarios finales (Pacheco et al., 2017). 
Además de fortalecer la conectividad interna del país, la RDNFO fue concebida para integrarse con redes internacionales mediante puntos de interconexión fronterizos. El proyecto contempla enlaces hacia Brasil (Iñapari), Bolivia (Desaguadero), Chile (Tacna) y Ecuador (Suyo), favoreciendo la integración regional y proporcionando rutas alternativas para el intercambio internacional de tráfico. Estas conexiones incrementan la resiliencia de la infraestructura nacional y contribuyen a mejorar la continuidad del servicio frente a posibles contingencias que afecten alguno de los enlaces principales (MTC, 2013).

\subsection{El modelo de operación mayorista}
\subsubsection{Concepto y marco regulatorio del <<Operador de Operadores>>}
\subsubsection{Estructura de la tarifa única y relación con operadores minoristas}
